\documentclass[11pt,a4paper]{article}
\usepackage[margin=1in]{geometry}
\usepackage{amsmath,amssymb,amsthm,booktabs,microtype,hyperref,url}
\hypersetup{colorlinks=true,linkcolor=blue,citecolor=blue,urlcolor=blue}

\newtheorem{theorem}{Theorem}[section]
\newtheorem{lemma}[theorem]{Lemma}
\newtheorem{proposition}[theorem]{Proposition}
\newtheorem{corollary}[theorem]{Corollary}
\theoremstyle{definition}
\newtheorem{definition}[theorem]{Definition}
\theoremstyle{remark}
\newtheorem{remark}[theorem]{Remark}

\newcommand{\F}{\mathbb{F}}
\newcommand{\Z}{\mathbb{Z}}

\title{Deletion-minor exclusions, heredity and telescoping-sum propagation\\
for Costas arrays of order 32}
\author{}
\date{}

\begin{document}
\maketitle

\begin{abstract}
Orders 32 and 33 are the smallest for which the existence of a Costas array is
unknown. We report four results that sharply delimit where an order-32 array
could lie and introduce a new propagation rule for distinct-difference search.

First, two census-complete exclusions: no order-32 Costas array contains an
order-29 Costas array as a three-row, three-column deletion minor, nor an
order-28 Costas array as a four-row, four-column deletion minor. The second
quantifies over all $115{,}087{,}822{,}400$ core placements of the complete
order-28 census. We show the two are logically independent.

Second, a measurement over the complete order-28 census showing that the Costas
property is almost never hereditary: $14.6\%$ of arrays contain an order-27
Costas array, $3.4\%$ an order-26 one, and none an order-25 one. Of the 112
order-27 sub-arrays found, 104 arise from corner deletion, which we explain: a
corner deletion induces a translation on the survivors while any interior
deletion induces a compression.

Third, a coordinate classification of the complete order-28 census. Screening in
exponential coordinates assigns every array a field and a degree using no
provenance information, and reproduces the published per-construction breakdown
$\{84,3,1,1\}$ exactly, class by class.

Fourth, a telescoping-sum propagation rule for distinct-difference constraints.
We prove it sound, formalise the proof in Lean~4 without Mathlib, and evaluate it
under a prospectively frozen, cost-controlled protocol: on the complete order-12
benchmark it detects $117{,}880$ of $296{,}656$ infeasible prefixes against
$77{,}074$ for the baseline, at $1.050\times$ the measured propagation work and
with no completable prefix rejected.

All catalogues, programs, machine-checked proofs and raw run logs are released
with checksums and a clean-room reproduction procedure.
\end{abstract}

\section{Introduction}

A \emph{Costas array} of order $n$ is a permutation $h$ of $\{0,\dots,n-1\}$ such
that the $\binom{n}{2}$ displacement vectors $(j-i,\,h_j-h_i)$ are pairwise
distinct. They arise in sonar and radar waveform design and have been studied as
combinatorial objects since Costas's 1965 report.

Two algebraic families, due to Welch and to Golomb, produce arrays at orders
determined by the prime powers nearby, and they remain the only generally
applicable constructions known. Exhaustive enumeration has reached order 29, at
a cost of $366.55$ CPU-years; orders 30 and 31 have never been enumerated,
although arrays are known at both. Orders 32 and 33 are the smallest at which
none is known.

We do not settle the existence question. We delimit it. Section~\ref{sec:ladder}
closes two extension channels completely; Section~\ref{sec:hered} measures how
much of the problem those channels ever covered; Section~\ref{sec:coord}
classifies the last completely enumerated census that is construction-explained;
and Section~\ref{sec:tm} contributes a propagation rule that is not specific to
Costas arrays and should apply to distinct-difference problems generally.

\paragraph{Reproducibility.} Every computation is released with source,
catalogues, raw logs and checksums, and every exhaustive claim is produced by a
fail-closed pipeline that refuses a verdict on incomplete coverage
(Section~\ref{sec:repro}).

\section{Preliminaries}

The dihedral group $D_4$ acts on the board; the Costas property is invariant, so
arrays come in equivalence classes of size 8, or 4 when symmetric.

Write $W_1(p,g,c)$ for the Welch array $h_i = g^{i+c} \bmod p - 1$ of order
$p-1$, and $G_2(q,\alpha,\beta)$ for the Golomb array over $\F_q$ of order $q-2$
with a dot at $(i,j)$ iff $\alpha^i+\beta^j=1$. Deleting corner dots from these
gives the variants conventionally labelled $W_2,W_3,G_3,G_4,G_5$.

\begin{definition}[Deletion minor]
An order-$m$ array $K$ is a \emph{deletion minor} of an order-$n$ array $A$ if
deleting $n-m$ columns of $A$, together with the rows they occupy, and
compressing the remainder yields $K$.
\end{definition}

\begin{definition}[Augmentation rung]
Fix $k$ and a Costas core $K$ of order $n-k$. Choose $k$ columns and $k$ rows of
the $n\times n$ board, place the core dots at the surviving columns and rows in
increasing order, and place $k$ further dots as a bijection between the chosen
columns and rows. Rung $k$ is the exhaustive search over all such configurations.
\end{definition}

Rung $k$ enumerates exactly the order-$n$ permutations possessing a given
order-$(n-k)$ Costas deletion minor, so an exhaustive rung with no solution is an
exclusion theorem.

\section{Catalogue generation and census verification}
\label{sec:cat}

The rung searches require catalogues of cores. We generate them rather than
import them, which makes the input auditable.

Seeding with $W_1(29)$, $W_1(31)$, $G_2(29)$, $G_2(31)$, $G_2(32)$ and closing
under corner deletion \emph{and corner addition}, then under $D_4$, produces 712
order-28 arrays in 89 dihedral orbits.

\begin{proposition}\label{prop:cat28}
The generated order-28 catalogue is the complete census.
\end{proposition}
\begin{proof}
The published enumeration is exhaustive with 712 arrays. The generated catalogue
contains 712 distinct permutations of $\{0,\dots,27\}$, each verified to have
Costas defect 0, closed under $D_4$ with 89 orbits. Two sets of equal finite
cardinality, one contained in the other, are equal.
\end{proof}

Corner \emph{addition} is essential and was initially omitted: $G_2(29)$ has
order 27 and reaches 28 only upwards. A downward-only closure produces 704 arrays
in 88 orbits, missing exactly one orbit, with representative
\[
(0,16,21,12,18,22,23,26,10,7,24,14,27,6,17,5,25,11,20,2,4,19,15,13,8,1,9,3).
\]

At order 29 the same closure produces 148 of the 164 census arrays. The
shortfall is exactly the two $RW_0$ equivalence classes produced by the
semi-empirical Rickard--Welch method, which appends a blank row, cyclically
shifts, and adjoins a column --- not a corner operation, so no corner closure
reaches it. We reconstruct the complete order-29 census from the 23 published
equivalence-class representatives, validating each independently as a Costas
permutation and reproducing every published figure: 164 arrays, 23 classes, 5
symmetric, per-method counts $4/64/24/24/32/16$.

\begin{proposition}[Order reachability]\label{prop:reach}
Closing every $W_1(p)$ for prime $p\le43$ and every $G_2(q)$ for prime power
$q\le43$ under corner deletion and addition reaches orders
$\{26,\dots,30\}\cup\{34,35,36\}$ and no others in $[20,34]$. In particular
$G_2(37)$ reaches 34 by one deletion and admits no second, so the families
conventionally labelled $G_4$ and $G_5$ are empty at $q=37$, and no corner-surgery
descendant of any such seed has order 31, 32 or 33.
\end{proposition}

\section{Two census-complete deletion-minor exclusions}
\label{sec:ladder}

All rungs use core-first pruning --- rejecting a placement as soon as the placed
core carries a repeated displacement, which is exact for defect 0 --- with
independent re-scoring of every low-defect configuration and a fail-closed
coverage manifest.

\begin{theorem}\label{thm:main}
No order-32 Costas array contains an order-29 Costas array as a three-row,
three-column deletion minor. No order-32 Costas array contains an order-28 Costas
array as a four-row, four-column deletion minor.
\end{theorem}

\begin{center}
\begin{tabular}{@{}lrrrr@{}}
\toprule
rung & cores & measured placements & Costas found & minimum defect \\
\midrule
$k=3$ (order 29) & 164 arrays, 23 classes & $565{,}836{,}800$ & \textbf{0} & 4 \\
$k=4$ (order 28) & 712 arrays, 89 classes & $115{,}087{,}822{,}400$ & \textbf{0} & --- \\
\bottomrule
\end{tabular}
\end{center}

The $k=4$ run was executed as 89 deterministic work units, one dihedral orbit
each; all report \texttt{EXHAUSTED}, and the fail-closed merge returns
\texttt{EXHAUSTED\_AND\_VERIFIED}. Because order 28 has no sporadic arrays,
``census-complete'' and ``all algebraic'' coincide there.

For comparison, rung $k=2$ over the 664 generated order-30 cores is exhaustive
relative to that catalogue --- order 30 has never been enumerated --- and yields
no order-32 array, with minimum defect 2 attained in exactly one dihedral orbit.
That orbit is the known near-solution, and rung $k=1$ identifies its parent: the
unique order-31 Costas orbit reachable by one-point augmentation of the algebraic
order-30 cores.

\begin{remark}[Independence]
Neither clause of Theorem~\ref{thm:main} implies the other. Deriving the $k=4$
clause from the $k=3$ clause would require every order-29 Costas array to contain
an order-28 Costas deletion minor, and conversely. Neither implication is
established, and Theorem~\ref{thm:hered} shows the required deletion property is
far from automatic. The two close different extension channels.
\end{remark}

\section{Heredity is corner-driven}
\label{sec:hered}

How much of the problem do such exclusions cover? With a complete census at order
28 this is directly measurable.

\begin{theorem}\label{thm:hered}
Over the complete order-28 census, the fraction of arrays possessing a Costas
deletion minor of order $28-k$ is $14.6\%$ for $k=1$, $3.4\%$ for $k=2$ and $0\%$
for $k=3$. Of the 112 order-27 sub-arrays found, 104 arise from corner deletion
and 8 are non-corner coincidences.
\end{theorem}

The explanation is immediate and is the reason the figure is so low: a corner
deletion removes a dot at a boundary of both ranges, so the compression acting on
the survivors is a translation, which preserves every displacement. Any interior
deletion induces a genuine compression, which does not. The measured corner rate
of $14.6\%$ is indistinguishable from the random-permutation rate $4/28=14.3\%$.

This bounds the reach of the entire constructed-and-augmented approach, and
correspondingly enlarges the residual class that no current structural theory
controls.

\section{Coordinate classification of the order-28 census}
\label{sec:coord}

The Welch relation $h_i = g^{i+c}-1$ is linear in the \emph{multiplicative}
parameter $u_i=g^i$, not in the additive index. Interpolating against additive
nodes therefore destroys exactly the structure it seeks: calibrated on known
algebraic arrays it returns maximal degree and full support, indistinguishable
from random permutations.

In exponential coordinates the families separate. Put $u_i=g^i\in\F_q$ against
$v_i=h_i$ (single-exponential, Welch), or $u_i=\alpha^{i+1}$ against
$v_i=\beta^{h_i+1}$ (dual-exponential, Golomb, where $\alpha^i+\beta^j=1$ gives
$v=1-u$). Offsets are free, since $u\mapsto\lambda u$ sends $a_j\mapsto
a_j\lambda^{-j}$ and $v\mapsto\mu v$ scales the interpolant.

\begin{theorem}\label{thm:coord}
Every array in the complete order-28 census has degree 1 in single-exponential
coordinates over $\F_{29}$, or in dual-exponential coordinates over $\F_{31}$ or
$\F_{2^5}$, except for a single dihedral orbit of 8 arrays generic in both. That
orbit is $G_2(29)$ of order 27 augmented by one corner dot.
\end{theorem}

\begin{center}
\begin{tabular}{@{}lrrll@{}}
\toprule
class & arrays & orbits & coordinates & degree \\
\midrule
Welch direct & 672 & 84 & single, $\F_{29}$ & 1 \\
Golomb, one deletion & 8 & 1 & dual, $\F_{31}$ & 1 \\
Golomb, two deletions & 24 & 3 & dual, $\F_{2^5}$ & 1 \\
augmented residue & 8 & 1 & generic in both & --- \\
\bottomrule
\end{tabular}
\end{center}

The screen uses no provenance information, yet the published enumeration's
breakdown into 84 $W_1$, 3 $G_4$, 1 $G_3$ and 1 $G_1$ classes gives the same
multiset $\{84,3,1,1\}$ class by class. Extension-field arithmetic is not
optional: $q-2=28$ needs $q=30$, not a prime power, so no intact Golomb array of
order 28 exists and every Golomb-type array there is a deletion from $\F_{31}$ or
$\F_{2^5}$.

The governing asymmetry follows: corner deletion is a translation and preserves
the degree-1 parameterisation, whereas adjoining a point off the curve forces the
interpolant to full degree $N-1$, since $P-L$ vanishes at the other $N-1$ nodes.

\section{Telescoping-sum propagation}
\label{sec:tm}

For a partial assignment and a lag $m$, the identity
\[
\sum_{i}\bigl(h_{i+m}-h_i\bigr)=\Bigl(\sum\text{last }m\Bigr)-\Bigl(\sum\text{first }m\Bigr)=:T_m
\]
holds for every sequence, and yields $T_m=T_{n-m}$ for free. It couples the
lag-$m$ difference multiset to the endpoint sums, which no per-lag
\textsc{AllDifferent} can see.

\begin{lemma}[Soundness]\label{lem:sound}
Let $k$ be the number of lag-$m$ pairs not wholly inside the prefix, $U$ the union
of their difference domains, $\det$ the sum of the determined differences, and
$R$ the remaining value set. Let $[\ell,h]$ bound $\det$ plus the sum of $k$
distinct elements of $U$, and $[\tau_-,\tau_+]$ bound $T_m$ by extremal subset
sums of $R$. If $[\ell,h]\cap[\tau_-,\tau_+]=\varnothing$ there is no Costas
completion.
\end{lemma}
\begin{proof}
For any completion $\sigma$ the identity above holds. Since $\sigma$ is Costas,
its lag-$m$ differences are pairwise distinct, so the $k$ undetermined ones are
$k$ distinct elements of $U$, whose sum lies between that of the $k$ smallest and
the $k$ largest. The unknown parts of the endpoint sums are subset sums of $R$ of
fixed size. Both constructions are outer approximations: they may admit sums no
completion realises, but exclude none that a completion realises. The two sides
being equal, the intervals meet.
\end{proof}

Lemma~\ref{lem:sound} and the surrounding development are formalised in Lean~4
(18 modules, no Mathlib, no \texttt{sorry}, no \texttt{native\_decide}),
including the extremal lemma that $k$ distinct elements of a finite integer set
have sum between those of its $k$ smallest and $k$ largest, and soundness of the
underlying domain filter.

\paragraph{Evaluation.} The protocol was frozen before the run. The benchmark is
the complete population of Costas-valid six-leading-column prefixes at order 12,
labelled by exhaustive completion: 304{,}288 valid, 296{,}656 infeasible,
7{,}632 completable, the split independently reproduced by a second completion
procedure. The baseline is relation-algebra closure united with position--value
domains under unary filtering, arc consistency, Hall matching and singleton
shaving.

\begin{theorem}\label{thm:tm}
On this population the baseline detects $77{,}074$ infeasible prefixes. Adding
the $T_m$ rule detects $117{,}880$, including $40{,}806$ additional, at $1.050$
times the baseline measured propagation work. Neither rejects any of the
$7{,}632$ completable prefixes.
\end{theorem}

The frozen gate required $\lceil 1.104167\times77{,}074\rceil = 85{,}103$
detections, zero false positives and cost ratio at most 2; all three are met, the
first by $32{,}777$. The result reproduces at five further complete populations
at orders 10 to 12.

\begin{remark}[Tree reduction is not speed-up]
Under a separately frozen protocol, adding $T_m$ to an exact solver removes
$26$--$32\%$ of the remaining search tree at orders 11 to 13 for about $4\%$
additional deterministic work, with solution sets verified identical by exact
byte comparison. It nevertheless runs $12$--$18\%$ slower in wall time, and both
propagation configurations are $53$--$63\times$ slower than a bare incremental
bitmask check that explores $5.5$--$8.1\times$ more nodes. The rule is a correct
and powerful deduction whose current implementation does not pay for itself in
seconds; the bottleneck is singleton shaving, not $T_m$.
\end{remark}

\section{Reproducibility}
\label{sec:repro}

Exhaustive negative results are only as good as the discipline that produced
them. Four conventions were adopted after failures in this work.

\emph{Assert the invariant before computing on the object.} Every tool validates
its input as a Costas permutation and halts otherwise.

\emph{Test the positive arm.} Detection power is measured on infeasible inputs;
soundness can only be measured on feasible ones. Two propagator defects arose
here --- a singleton-shaving routine that silently reinitialised its domains, and
a relation filter that emptied every relation --- and the second was invisible to
detection counts and immediate under a completable-arm test. Both were
deterministic and internally consistent, and both would have passed a certificate
of their own output.

\emph{Compute the null analytically when sampling cannot resolve it.} Three
analyses in this work drew conclusions from samples far below the resolution of
the quantity estimated.

\emph{Fail closed on coverage.} The merge tool refuses an exhaustive verdict
unless every unit reports exhaustion, all agree on parameters and catalogue hash,
coverage is exact with no gap or duplicate, and every reported configuration
re-scores from scratch.

\section{Conclusion}

Order 32 remains open. What has changed is the size of the space in which an
array could lie: two deletion-minor channels are closed over complete censuses, a
third is closed relative to the generated order-30 catalogue, no corner-surgery
descendant of any Welch or Golomb seed has order 31, 32 or 33, and the order-28
census is completely explained by coordinate system. Against this,
Theorem~\ref{thm:hered} shows that deletion minors were never a large fraction of
the problem, so the residual class is correspondingly larger than the closures
suggest --- and no current structural theory controls it.

The telescoping-sum rule is the component we expect to outlast the specific
question. It is sound, machine-checked, and couples information across
distance classes, which is where additional power must come from once per-class
filtering saturates.

\end{document}
